\documentclass{article}
\usepackage[letterpaper, margin=1in]{geometry}

%%%%%i added this
\newtheorem{theorem}{Theorem}{}
\newtheorem{corollary}{Corollary}{}
\newtheorem{remark}{Remark}{}
\usepackage[table]{xcolor}
\usepackage{type1cm}
\usepackage{lettrine}
\usepackage{xcolor}
\usepackage{hyperref}
\usepackage{subfigure}
\usepackage{multirow}
\usepackage[normalem]{ulem}
\useunder{\uline}{\ul}{}
\usepackage{amsmath,amsfonts}
\usepackage[left=0.9cm,right=0.9cm,top=1.9cm]{geometry}
\usepackage{array}
%\usepackage{algorithmic}
%\usepackage[ruled,linesnumbered]{algorithm2e}
\usepackage{amsmath,amssymb,amsfonts, bm}
\usepackage{algorithmic}
\usepackage[ruled,linesnumbered]{algorithm2e}

%%%%%%%


\usepackage[utf8]{inputenc} % make more characters printable when copy-pasted here
\usepackage{xcolor} % to facilitate the myNotes and myTodo commands.
\usepackage{amsmath}
\usepackage{graphicx}
\usepackage{subcaption}
\usepackage{wrapfig}
\usepackage{titling}
\usepackage{soul}
\usepackage{hyperref}
\usepackage{booktabs}
\usepackage{cite}
\usepackage[linesnumbered,ruled,vlined]{algorithm2e}%
\usepackage[linesnumbered,ruled,vlined]{algorithm2e}%

\renewcommand\maketitlehooka{\null\mbox{}\vfill}
\renewcommand\maketitlehookd{\vfill\null}
\usepackage[article]{ragged2e}
% Created by Jack Baker, last modified 11/3/2020

% initialize counters for reviewer number and comment number
\newcounter{reviewer}
\setcounter{reviewer}{0}
\newcounter{point}[reviewer]
\setcounter{point}{0}
\setcounter{figure}{2}
\renewcommand{\thetable}{\Roman{table}}
\setcounter{table}{0}
\usepackage{blindtext}
\usepackage{parskip}
\usepackage{caption} 
\captionsetup[table]{skip=10pt}

\usepackage{multirow}
\usepackage{pifont}% http://ctan.org/pkg/pifont
\newcommand{\cmark}{\ding{51}}%
\newcommand{\xmark}{\ding{55}}%
% Set the format of reviewer comment numbers and responses.
\renewcommand{\thepoint}{Review\,\thereviewer.\arabic{point}} 
\newcommand{\point}[1]{\refstepcounter{point}  \bigskip \hrule \medskip \noindent \textit{{\fontseries{b}\selectfont \thepoint } #1}\par}
\newcommand{\reply}{\medskip \noindent \textbf{Reply}:\ }   

% Set the format for reviewer headings
\newcommand{\reviewersection}{\stepcounter{reviewer}
                  \section*{To Reviewer \#\thereviewer}}
                  
% Set the format for reviewer headings
\newcommand{\editorsection}{\section*{To Associate Editor}}

% Set the format for notes on in-progress responses
\newcommand\myNotes[1]{\textcolor{red}{#1}}
\newcommand\myTodo[1]{\textcolor{blue}{#1}}

%Table Space
\newcommand\T{\rule{0pt}{2.5ex}}       % Top strut
\newcommand\B{\rule[-1.3ex]{0pt}{0pt}} % Bottom strut


\hypersetup{
    colorlinks,
    linkcolor={blue!80!black},
    citecolor={blue!80!black},
    urlcolor={blue!80!black}
}




%%%%%%%%%%%%%%%%%%%%%%%%%%%%%%%%%%%%%%%%%%%%%%%%%%%%%%%%
\title{OrBIT: Orchestrated Hybrid-Blockchain Twins for Verifiable Privacy-Preserving Provenance \\ (Paper Number: IOT-D-26-03176 )}
\author{}
\begin{document}

\maketitle

\vspace{10 pt}

% \noindent
% \textit{[Jack Baker comments: This is a simple Latex template for listing your paper review comments and responses. 
% It could easily be done in Microsoft Word as well, but I prefer Latex if the original paper is in Latex, because it is easier to copy-and-paste passages from the paper into the replies.
% The template has a few simple commands (myNotes and myTodo) that I like to use to leave colored text passages of in-progress items, but otherwise is a bare-bones template.]}


% \vspace{20 pt}

% \noindent \textit{[A place for general comments regarding the comments and responses, plus a than-you to the reviewers.]}

% The reviewer's comments have been listed in order below in italic text, followed by our responses in normal text. 
% \pagewiselinenumbers
\newpage
\editorsection
\justifying{Thank you for your thoughtful remarks. According to your letter, our paper (IOT-D-26-03176) goes through mandatory major revisions based on the comments of reviewers 1 and 2. Regarding the corrections and suggestions from the reviewer, we have modified the manuscripts.  Below are the reviewer's remarks in italic text, followed by our reply in regular style, and the manuscript modifications are highlighted in \textcolor{blue}{blue}. Table \ref{edited} summarizes major differences between the old version and the enhanced version affected based on the review comments. 

 \begin{table}[!thbp]
    \caption{Summary of major differences between the previous manuscript and the Revised version.}
    \label{edited}
    \centering
    \begin{tabular}{l|c|c}
    \hline\hline
    Paper Changes   & Old Manuscript& Revised Manuscript\T\B\\ \hline\hline
    Number of pages                         & 18  & 19  \T\B\\\hline
        Number of Figures                      & 7 &  10 \T\B\\\hline
        Number of Tables                            & 7 & 4     \T\B\\\hline

        Number of References                               & 17 &   24 \T\B\\ \hline\hline
    \multicolumn{2}{l}{\parbox{3.3in}{}}
    \end{tabular}
    \vspace{-4mm}
    \end{table}

\newpage
\reviewersection

We thank the reviewer for identifying limitations in the original experimental
design and for requesting stronger evidence on distribution, relay fault
tolerance, and comparison baselines. These comments motivated the new cloud
deployment and most of the additional experiments in the revision.

\point{The entire empirical evaluation, including Hyperledger Fabric consensus,
the local Ethereum RPC, and the orchestrator, was conducted on a single 4-core
machine running development WSGI servers. This completely masks the actual
network latency, multi-organization endorsement overhead, and Byzantine fault
conditions that define real-world distributed ledger deployments, rendering the
throughput and latency claims practically meaningless. Please deploy and
benchmark the architecture across a geographically distributed, multi-node
cloud environment with realistic network delays and production-grade
application servers.}

\reply{We agree that the original single-host deployment could not support
claims about distributed-ledger latency. We therefore rebuilt the evaluation on
four Azure virtual machines spanning Korea Central and Japan East. The revised
network separates the orchestrator/client, two Fabric peer organisations and
their CouchDB instances, and the ordering/public-ledger services. Fabric uses an
\texttt{AND(Org1,Org2)} endorsement policy, and the Flask services run behind
Gunicorn. We also retained the original host only for a controlled
Werkzeug-versus-Gunicorn experiment, which showed that changing the WSGI server
alone did not materially change throughput; Fabric ordering and commit remained
the dominant cost.

The new experiments include controlled 0--100\,ms added round-trip delay,
\texttt{MaxMessageCount} and \texttt{BatchTimeout} sweeps, separate endorsement
and commit-stage measurements, and five repetitions per configuration. The
base distributed run exhibited acknowledgement floors of 2637.5--2664.9\,ms
and a transition at $c=10$, while reducing
\texttt{MaxMessageCount} to 5 moved the detected transition to $c=5$. A stage
probe measured 319.58\,ms for endorsement/submission and 1742.67\,ms for
ordering/commit.

We could not responsibly claim Byzantine-fault evaluation. The deployed Fabric
v2.5 network uses a single Raft consenter, whereas a Byzantine-ordering
experiment requires migration to Fabric v3 and at least four orderers. We now
state this limitation explicitly throughout the results, limitations, and
conclusion. Thus, the geographic, multi-organisation, delay, and
production-server portions of the request are fully addressed; Byzantine
ordering is transparently reported as not run.}

\begin{quote}
\textbf{Changes in the manuscript (Experimental Setup, Ledger and Network
Configuration, Results and Discussion, Figures~4--9):}

\textcolor{blue}{``The revised OrBIT evaluation uses four Azure virtual machines
distributed across Korea Central and Japan East. The client/orchestrator host,
an Org1 peer with CouchDB, an Org2 peer with CouchDB, and the ordering/public-ledger
host are separated at VM and regional boundaries. Hyperledger Fabric~v2.5
enforces \texttt{AND(Org1,Org2)} endorsement, so every accepted update crosses
both administrative organisations and regions before ordering and validation.''}

\textcolor{blue}{``The deployed ordering service has one Raft consenter and
therefore tolerates no orderer failure. Fabric~3 Byzantine ordering would
require a full network migration and at least four suitably provisioned
orderers; it was not run under the available quota.''}
\end{quote}

\point{The architecture relies heavily on a centralized ``Orchestrator'' and a
single-threaded ``Relay Worker'' to guarantee exactly-once delivery between the
private and public ledgers. Centralizing the cross-chain synchronization
eliminates the decentralized trust assumption of the digital twin ecosystem,
introducing a critical single point of failure that compromises the promised
end-to-end immutability. Please replace the centralized relay with a
decentralized oracle network or a cryptographically secure cross-chain bridge
protocol, and formally evaluate its fault tolerance.}

\reply{We agree with the underlying concern and have removed the unsupported
implication that a single relay provides decentralised or trustless bridge
security. We did not replace OrBIT with an external oracle network because that
would change the system's trust, validator, incentive, and finality models rather
than test the contribution under review. Instead, we made the actual permissioned
relay design fault tolerant and its trust boundary explicit.

The revised implementation supports $k\in\{1,2,3\}$ independently restartable
relay processes with separate Ethereum accounts. We compare atomic lease-based
coordination against an uncoordinated race mode and preserve the on-chain
identifier-uniqueness check as the final exactly-once safety boundary. We added
fault injection at mid-poll, after broadcast, after receipt, and by
\texttt{SIGKILL}. In every lease-coordinated case, anchors equalled records,
duplicate anchors and reverted transactions were zero, and delivery recovered
after lease expiry. The independent auditor detected all commitment-tampering
and withholding trials. A separate crash-window experiment created 150 private
ledger orphans; reconciliation found and repaired all 150, after which all 150
anchors passed audit.

These tests establish crash-fault tolerance and externally detectable
cross-ledger integrity. They do not establish Byzantine relay consensus or a
trustless bridge, and the revised manuscript now says so explicitly. A
decentralised oracle or light-client bridge can consume the same commitment and
outbox interface and is identified as future comparative work.}

\begin{quote}
\textbf{Changes in the manuscript (Related Works, Orchestrator Integration
Layer, Algorithm~1, Fault-Tolerance Evaluation, Limitations, and Conclusion):}

\textcolor{blue}{``Exactly-once effects arise from three mechanisms acting
together. The outbox key prevents duplicate enqueueing; relay leases allocate
pending work without relying on one continuously available process; and the
on-chain uniqueness check closes the crash window between transaction broadcast
and receipt persistence.''}

\textcolor{blue}{``The evidence above supports crash-fault-tolerant relay
coordination, not Byzantine consensus or a trustless cross-chain bridge.''}
\end{quote}

\point{The baseline comparisons are glaringly insular, evaluating OrBIT's
hybrid routing exclusively against its own ``Fabric-only'' ablation rather than
competing state-of-the-art hybrid twin architectures or federated cross-chain
protocols. While the theoretical integration of immutable public anchoring with
private state management is sound, treating a centralized, single-threaded
Python relay as a cross-chain bridge represents a regression relative to
decentralized oracle networks or trustless relay models that are standard in
contemporary blockchain literature. Check them for improvement:
\nolinkurl{10.1109/TCE.2026.3660345},
\nolinkurl{10.1038/s41598-026-46244-z}, and
\nolinkurl{10.1109/TCE.2026.3702497}.}

\reply{We expanded both the literature analysis and the experimental baselines.
The three suggested papers are now cited and discussed individually. They address
federated predictive digital twins, privacy-preserving medical-image learning,
and neural-cryptographic health-data exchange. They are relevant neighbouring
systems, but they do not expose the same dependent variables as OrBIT's
permissioned-to-public commit protocol; using their prediction accuracy or model
cost as a throughput baseline would therefore be methodologically invalid.

To make an experimentally controlled comparison, we added three protocol modes
with identical payloads and ledger topology: \texttt{FABRIC\_ONLY},
\texttt{SYNC\_ANCHOR}, and OrBIT's \texttt{ASYNC\_OUTBOX}. We also compare
lease-coordinated relays with uncoordinated race relays. At $c=32$, synchronous
anchoring reached 2.73\,TPS, whereas asynchronous and Fabric-only modes reached
9.05 and 9.09\,TPS, respectively. The manuscript now distinguishes this
permissioned, auditable delivery mechanism from oracle or trustless-bridge
consensus and avoids presenting the relay as a decentralised bridge.}

\begin{quote}
\textbf{Changes in the manuscript (Related Works and Baseline Comparisons):}

\textcolor{blue}{``The three additional works suggested during review occupy
adjacent rather than directly competing layers. CardioTwin-XAI builds a
federated predictive digital twin for cardiovascular care; FedLiverNet protects
collaborative liver-image learning; and NeuroShield-Diabetes applies neural
cryptography to health-data exchange. They provide application- or model-level
privacy mechanisms but do not specify a permissioned-to-public cross-chain
commit protocol, relay coordination rule, or comparable end-to-end ledger
workload.''}
\end{quote}

\reviewersection

We thank the reviewer for recognising the relevance and organisation of the
work and for the concrete suggestions on motivation, algorithms, presentation,
language, and recent literature. All five requested revisions have been made.

\point{Problem of paper and motivation is not clear in introduction.}

\reply{We rewrote the opening of the Introduction to state the engineering
problem before introducing the solution. It now identifies three coupled risks:
disclosure of regulated fields, duplicate or missing cross-ledger effects after
retries or crashes, and public-chain delay on the ingestion path. We also
rewrote the contribution list so that every contribution identifies both a
mechanism and a verifiable outcome.}

\begin{quote}
\textbf{Change in the manuscript (Introduction):}

\textcolor{blue}{``The central problem is not merely where to store digital-twin
records, but how to coordinate one private and one public ledger without
(i) disclosing regulated fields, (ii) producing duplicate or missing
cross-ledger effects after retries and crashes, or (iii) placing public-chain
delay on the ingestion critical path.''}
\end{quote}

\point{Algorithm is discussed but no pseudo code is added for ease of readers.}

\reply{We added three numbered pseudocode algorithms covering the paper's core
procedures: (1) lease-coordinated relay delivery and reconciliation,
(2) deterministic policy-driven routing and commit, and (3) digital-twin
version commit and rollback. The surrounding text now refers to each algorithm
at the point where its inputs, safety condition, and output are introduced.}

\begin{quote}
\textbf{Changes in the manuscript (Algorithms~1--3):}

\textcolor{blue}{Algorithm~1 specifies lease acquisition, on-chain uniqueness,
receipt recovery, and reconciliation; Algorithm~2 specifies deterministic
field classification, fail-closed routing, Fabric commit, and outbox insertion;
Algorithm~3 specifies delta-or-snapshot version storage, checkpoint-bounded
reconstruction, digest verification, and rollback.}
\end{quote}

\point{Figure 1 colors are light and not clear so add increase size or enlarge
for clear view.}

\reply{Figure~1 is now a two-column \texttt{figure*} rendered at
0.80\,\texttt{textwidth}, rather than a single-column figure. We also moved its
in-text introduction before the figure and expanded the caption so that the
data and control directions remain interpretable at print scale. During the
same figure audit, we introduced and interpreted every experimental figure in
the body text before its corresponding float.}

\begin{quote}
\textbf{Change in the manuscript (Introduction, Figure~1):}

\textcolor{blue}{``Figure~1 summarises the resulting closed loop: IoT
observations flow from physical assets through edge processing and
hybrid-ledger governance to the digital twin, while verified decisions and
lifecycle updates flow back to the physical system.''}
\end{quote}

\point{Paper contains few grammar mistakes which will be cooperated in final
version.}

\reply{We edited the complete manuscript for grammar, punctuation, article use,
terminology, and consistency. We also standardised ``Figure'' references,
hyphenation such as ``cross-ledger'' and ``multi-region'', symbol definitions,
and the distinction between acknowledgement latency and completion latency.
Captions were revised to state what is measured and how each panel should be
interpreted.}

\point{Only few references are added in paper, but more than 40 references so
to attract readers add few latest references related to this paper, which is
mentioned below: Jumani et al., ``Blockchain Security and Privacy: Threats,
Solutions, and Future Directions''; Li et al., ``Blockchain applications for
Internet of Things (IoT): A review''; Laghari et al., ``Internet of Things for
gaming: A review''; and Laghari et al., ``Internet of multimedia things (IoMT):
A review.''}

\reply{We added and discussed all four suggested references in the Related
Works section. They are used to strengthen the manuscript's treatment of
blockchain security and privacy, blockchain--IoT integration, heterogeneous
application domains, and latency- and privacy-sensitive multimedia streams.
We also added the three recent papers suggested by Reviewer~1, for seven new
review-suggested references in total.}

\begin{quote}
\textbf{Change in the manuscript (Related Works):}

\textcolor{blue}{``Recent surveys sharpen the security and deployment context
of this problem. Laghari et al. catalogue blockchain--IoT applications and
identify interoperability, scalability, and multi-organisation security as
persistent constraints; Jumani et al. similarly organise blockchain threats,
privacy mechanisms, and quality-of-service trade-offs. IoT reviews in gaming
and multimedia environments show that heterogeneous sensors generate
latency-sensitive and privacy-sensitive streams beyond conventional industrial
telemetry.''}
\end{quote}

\section*{Closing Statement}

We again thank the reviewers for comments that materially improved the paper.
The revision now provides a geographically distributed evaluation, explicit
fault and trust boundaries, protocol-comparable baselines, reproducible
pseudocode, clearer figures, and a point-by-point account of both completed and
unrun experiments. We have deliberately avoided claiming Byzantine or trustless
bridge guarantees that were not measured.

\end{document}
